% La SVD: leer la forma de las curvas de nivel (partición por clase: apunte/scripts/reorganizar_por_clase.py)
\chapter{La SVD: leer la forma de las curvas de nivel}\label{cap:svd}

\section{La idea}


% --- 09-svd-nivel1 (redactado 2026-10-18) ---

\subsection{Una base para la entrada, otra para la salida}\label{sec:svd-idea}

\paragraph{La pregunta que cambia.} En el capítulo de álgebra buscamos las direcciones que
una transformación deja caer sobre sí mismas, los autovectores, y con ellas escribimos la
matriz como diagonal: \(A=PDP^{-1}\) (sección~\ref{sec:cambio-base}). La idea es buena y
tiene tres defectos que ya vimos. Una rotación no deja ninguna recta real fija; un corte deja
una sola; una matriz rectangular, que lleva el plano a \(\R^3\) o \(\R^3\) al plano, ni
siquiera puede preguntarse qué vectores caen sobre sí mismos, porque la entrada y la salida
son espacios distintos. La descomposición en valores singulares (SVD, por sus siglas en
inglés) cambia la pregunta: en lugar de una sola base que sirva a la vez para la entrada y
para la salida, buscamos \emph{dos} bases ortonormales, una en cada espacio, elegidas de
manera que la transformación lleve la primera sobre la segunda estirando cada eje por
separado. En esas dos grillas cualquier matriz, cuadrada o rectangular, con autovalores
reales o sin ellos, es una matriz diagonal con entradas no negativas:
\[
  A=U\,\Sigma\,V^T,\qquad
  V:\ \text{base de la entrada},\quad
  \Sigma:\ \text{estirar cada eje},\quad
  U:\ \text{base de la salida}.
\]
Lo que se gana es generalidad; lo que se pierde es la posibilidad de componer la
transformación consigo misma en esa base, porque la grilla de llegada no es la de partida.

\paragraph{Hecho geométrico: lo redondo se vuelve elipse.} La forma más directa de ver las
dos bases es mirar qué hace la transformación con todas las direcciones a la vez.

\begin{itemize}
  \item En \(\R^2\), la circunferencia unidad va a parar a una elipse (o, si la matriz
  aplasta, a un segmento). Los largos de los semiejes de esa elipse son los \emph{valores
  singulares} \(\sigma_1\ge\sigma_2\ge0\); las direcciones de los semiejes son las columnas
  de \(U\); las dos direcciones perpendiculares de la circunferencia que terminan sobre esos
  semiejes son las columnas de \(V\). Toda elipse tiene un par de ejes perpendiculares, y
  por eso la SVD existe siempre.
  \item En \(\R^3\), la esfera unidad va a parar a un elipsoide: tres semiejes
  perpendiculares, tres valores singulares, tres direcciones de llegada en \(U\) y tres de
  partida en \(V\). Si un \(\sigma\) es cero, el elipsoide se aplastó a una elipse plana, y
  la dirección de partida correspondiente es el núcleo.
  \item En una matriz de \(m\times n\), la esfera de \(\R^n\) va a parar a un elipsoide,
  posiblemente aplastado, dentro de \(\R^m\). La cantidad de semiejes no nulos es el rango.
  Las direcciones de \(\R^n\) que se aplastan forman el núcleo y las direcciones de
  \(\R^m\) que la elipse no alcanza forman el espacio inaccesible: los cuatro subespacios
  de la sección~\ref{sec:pseudoinversa} salen de la misma figura.
\end{itemize}

Leída de derecha a izquierda, la descomposición dice que toda transformación lineal es un
giro (o una reflexión) de la entrada, seguido de estiramientos sobre los ejes, seguido de un
giro (o una reflexión) en la salida. Los giros no cambian largos; todo lo que la
transformación hace con las distancias está concentrado en los números de \(\Sigma\).

\paragraph{Por qué el profesor la llama ``base privilegiada''.} Para el profesor, el punto
de partida de la SVD es encontrar una ``base privilegiada'' de la transformación (resumen
2025, clase 7). La palabra tiene peso, porque el mismo curso insiste en que, en abstracto,
``no existe una base vectorial privilegiada de \(\R^n\), y lo único que define a la base que
llamamos canónica es que sobre ella tenemos definida la noción de distancia a través del
producto interno'' (resumen 2025, clase 6). Y lo repite al hablar de cambio de base: ``En un
espacio vectorial donde la magnitud de un vector no sea algo que esté definido o que
importe, la base canónica no tiene nada especial'' (resumen 2025, clase 8).

Ahí está la clave. La SVD pide bases \emph{ortonormales}: vectores de largo uno y
perpendiculares entre sí. Largo y perpendicularidad no los da la transformación sino el
producto interno de cada espacio. Sin producto interno, una circunferencia no se distingue
de una elipse y no hay semiejes que medir; con producto interno, la elipse imagen tiene ejes
bien definidos y esos ejes eligen las bases. Por eso la base que la SVD privilegia no es
privilegiada por la matriz sola, sino por la matriz junto con las dos reglas de medir, la de
la entrada y la de la salida (sección~\ref{sec:puntos-formas-grillas}). Si cambiamos la
regla, cambian los valores singulares. Un ejemplo que vamos a usar con el terreno: la
transformación que lleva un desplazamiento de \((\Delta\text{columnas},\Delta\text{filas})\)
de una grilla geográfica a metros sobre el suelo es, en el recorte de la Precordillera,
\(\operatorname{diag}(26{,}19;\ 30{,}80)\), porque a \(32^\circ\) de latitud un segundo de longitud mide
aproximadamente \(\cos\varphi\approx0{,}85\) veces lo que mide un segundo de latitud. Para quien cuenta
celdas los dos ejes son iguales; para quien mide en metros, la circunferencia de radio una
celda es una elipse de semiejes \(30{,}80\) y \(26{,}19\) m.

\paragraph{Una herramienta, varios problemas de la materia.} La SVD reúne en una sola
cuenta preguntas que en el programa aparecen separadas. Cada una se contesta mirando una
parte distinta de \(U\Sigma V^T\):

{\small
\begin{center}
\begin{tabularx}{.94\textwidth}{>{\raggedright\arraybackslash}p{.42\textwidth} Y}
\toprule
Pregunta de campo & Qué parte de la SVD la contesta\\
\midrule
¿Cuántas de mis ecuaciones son de verdad independientes? ¿Un error chico en los datos se
vuelve un error grande en el resultado? &
cuántos \(\sigma_i\) son no nulos (el rango) y el cociente \(\sigma_1/\sigma_n\) (el
número de condición)\\
Tengo más observaciones que incógnitas y el sistema no cierra (una poligonal, un ajuste) &
la pseudoinversa \(A^{+}=V\Sigma^{+}U^T\): invertir los \(\sigma_i\) no nulos y descartar
lo que cae en el espacio inaccesible\\
Tengo los mismos puntos medidos en dos sistemas locales y busco el giro que mejor los
superpone (Procrustes) &
las dos bases \(U\) y \(V\) de una matriz chica armada con los datos; los \(\sigma_i\) no
intervienen en el giro\\
Ajusto una recta o un plano a una nube con error en todas las coordenadas (regresión total) &
el vector singular del \emph{menor} \(\sigma\): la dirección en que la nube es más delgada\\
¿Qué parte de un relieve es forma grande y qué parte es detalle? ¿Con qué precisión dibujo
sus curvas de nivel? &
los primeros términos \(\sigma_iu_iv_i^T\) de la matriz de alturas y la suma de los
\(\sigma_i^2\) que se descartan (sección~\ref{sec:svd-dem})\\
\bottomrule
\end{tabularx}
\end{center}
}

Las cuatro primeras filas se resuelven en las secciones~\ref{sec:pseudoinversa}
y~\ref{cap:cuatro-problemas}. La última es la que cierra el capítulo: un modelo digital de
elevación es una matriz de alturas, y su SVD separa el relieve en capas ordenadas de mayor
a menor importancia, a partir de las cuales trazamos curvas de nivel tan suaves o tan
detalladas como pida la escala del plano.

\linkUnidad{09_cambio_de_base_svd}{Simulación interactiva: Unidad 13 — Cambio de base, SVD e inversas}

\section{Construcciones, cuentas y ejemplos}


Reunimos tres piezas: la SVD, las inversas con los cuatro subespacios y la pseudoinversa, y la guía de los cuatro problemas. No se dictaron en 2C2026; las armamos a partir del resumen 2025 del profesor y de su cuaderno \file{LaSVD.wxmx}, y quedan pendientes de validar.

\subsection{Descomposición en valores singulares (SVD)}\label{sec:svd}

{\small\vspace{1mm}\linkUnidad{09_cambio_de_base_svd}{Simulación interactiva: Unidad 13 — Cambio de base, SVD e inversas}\par}
{\small\vspace{1mm}\linkAnimacion{04_svd_circulo_a_elipse}{Animación: SVD, del círculo a la elipse}\par}


\begin{atencion}[Anticipado]
La SVD no se dictó en 2C2026 (clase 7 del esqueleto; anunciada el 30/9 y el 14/10). Esta
sección se armó antes de que la cátedra la diera, a partir del resumen 2025 del profesor,
de su cuaderno \file{LaSVD.wxmx} y de bibliografía (Strang; Calvetti y Somersalo).
Pendiente de validar.
\end{atencion}

\paragraph{Cambiar la pregunta.} Diagonalizar pregunta qué direcciones caen sobre sí
mismas. Una rotación de \(90^\circ\) no deja ninguna, y el polinomio característico lo dice
con raíces complejas (sección~\ref{sec:autovalores}). El profesor planteó en 2025 que
``podría existir una forma de descomponer la transformación que no tuviera estos
problemas'', y la enunció así: ``siempre existe una base ortogonal del espacio de entrada
de una transformación que se conserva ortogonal a la salida''. La pregunta nueva es:
¿qué direcciones \emph{perpendiculares} de la entrada siguen perpendiculares a la salida?
Esa pregunta siempre tiene respuesta real, porque se construye a partir del producto
interno y no de direcciones fijas.

\paragraph{Hecho geométrico: la circunferencia se vuelve elipse.} Una matriz \(A\) de
\(2\times2\) lleva la circunferencia unidad a una elipse. Los semiejes de la elipse miden
\(\sigma_1\ge\sigma_2\ge0\) (valores singulares); sus direcciones \(u_1,u_2\) son una base
ortonormal de la salida; las direcciones de la circunferencia que van a parar a esos
semiejes, \(v_1,v_2\), son una base ortonormal de la entrada:
\[
  A\,v_i=\sigma_i\,u_i,\qquad
  \boxed{\ A=U\,\Sigma\,V^T\ },\qquad
  U=[u_1\ u_2],\ \ \Sigma=\operatorname{diag}(\sigma_1,\sigma_2),\ \ V=[v_1\ v_2].
\]
Leído de derecha a izquierda: \(V^T\) gira (o refleja) hasta poner \(v_1,v_2\) sobre los
ejes, \(\Sigma\) estira los ejes, \(U\) gira el resultado hasta las direcciones
\(u_1,u_2\). Tres movimientos que ya se entienden por separado. Como \(\Sigma\) es diagonal
con entradas no negativas, no hay giro escondido que arruine la cuenta: los giros se los
llevaron \(U\) y \(V\). Además \(\sigma_1\sigma_2=|\det A|\) (el área de la elipse es la de
la circunferencia por el producto de los semiejes), y \(\sigma_k=0\) significa que la
elipse se aplastó: \(v_k\) está en el núcleo y \(u_k\) queda fuera de la imagen.

Escalera. En \(\R\): \(x\mapsto ax\) se escribe \(a=(\pm1)\cdot|a|\cdot1\); el valor
singular es \(|a|\) y el signo quedó en la ``base de salida''. En \(\R^3\): esfera que se
vuelve elipsoide, tres semiejes. Si la matriz es \(3\times2\), como la jacobiana de una
superficie parametrizada, la circunferencia del plano \((u,v)\) va a una elipse que vive en
el plano tangente dentro de \(\R^3\) y el tercer vector \(u_3\) apunta a la normal, la
dirección inaccesible: \(U\) es \(3\times3\), \(V\) es \(2\times2\), \(\Sigma\) es
\(3\times2\). Es el ``cambio de base rectangular'' que el profesor anunció el 30/9.

\paragraph{Ejemplo a mano con enteros.} \(A=\begin{pmatrix}3&0\\4&5\end{pmatrix}\). La
base de entrada sale de \(A^TA\), que es simétrica:
\[
  A^TA=\begin{pmatrix}3&4\\0&5\end{pmatrix}\begin{pmatrix}3&0\\4&5\end{pmatrix}
  =\begin{pmatrix}25&20\\20&25\end{pmatrix},\qquad
  (25-\mu)^2-400=0\ \Rightarrow\ \mu=45,\ 5 .
\]
Autovectores \((1,1)\) y \((1,-1)\); por lo tanto \(\sigma_1=\sqrt{45}=3\sqrt5\approx6{,}71\),
\(\sigma_2=\sqrt5\approx2{,}24\), \(v_1=\tfrac1{\sqrt2}(1,1)\), \(v_2=\tfrac1{\sqrt2}(1,-1)\).
La base de salida se obtiene transformando, \(u_i=Av_i/\sigma_i\):
\[
  A\begin{pmatrix}1\\1\end{pmatrix}=\begin{pmatrix}3\\9\end{pmatrix}=3\begin{pmatrix}1\\3\end{pmatrix}
  \ \Rightarrow\ u_1=\tfrac1{\sqrt{10}}(1,3),\qquad
  A\begin{pmatrix}1\\-1\end{pmatrix}=\begin{pmatrix}3\\-1\end{pmatrix}
  \ \Rightarrow\ u_2=\tfrac1{\sqrt{10}}(3,-1).
\]
Chequeos: \(\|(3,9)\|/\|(1,1)\|=\sqrt{90}/\sqrt2=\sqrt{45}=\sigma_1\) y
\(\|(3,-1)\|/\|(1,-1)\|=\sqrt{10}/\sqrt2=\sqrt5=\sigma_2\); \(u_1\cdot u_2=(3-3)/10=0\);
\(\sigma_1\sigma_2=15=\det A\). La base ortogonal \(\{(1,1),(1,-1)\}\) de la entrada fue a
parar a la base ortogonal \(\{(1,3),(3,-1)\}\) de la salida.

\begin{center}
\begin{tikzpicture}[scale=.42]
  % entrada
  \begin{scope}
    \draw[axisline] (-3,0)--(3,0); \draw[axisline] (0,-3)--(0,3);
    \draw[mainline] (0,0) circle (2);
    \draw[rust,line width=1pt,-{Latex[length=2.2mm]}] (0,0)--(1.414,1.414)
      node[above right,font=\small,labelback] {\(v_1\)};
    \draw[tealdark,line width=1pt,-{Latex[length=2.2mm]}] (0,0)--(1.414,-1.414)
      node[below right,font=\small,labelback] {\(v_2\)};
    \node[text=muted,font=\small] at (0,-3.8) {entrada: circunferencia unidad};
  \end{scope}
  \draw[-{Latex[length=2.6mm]},line width=.9pt,muted] (4.2,0)--(6.8,0) node[midway,above,font=\small] {\(A\)};
  % salida
  \begin{scope}[xshift=14cm]
    \draw[axisline] (-7.5,0)--(7.5,0); \draw[axisline] (0,-7.5)--(0,7.5);
    \draw[mainline,rotate=71.565] (0,0) ellipse (6.708 and 2.236);
    \draw[rust,line width=1pt,-{Latex[length=2.2mm]}] (0,0)--(2.121,6.364)
      node[right,font=\small,labelback] {\(\sigma_1u_1=(3,9)/\sqrt2\)};
    \draw[tealdark,line width=1pt,-{Latex[length=2.2mm]}] (0,0)--(2.121,-0.707)
      node[below right,font=\small,labelback] {\(\sigma_2u_2=(3,-1)/\sqrt2\)};
    \node[text=muted,font=\small] at (0,-8.3) {salida: elipse de semiejes \(3\sqrt5\) y \(\sqrt5\)};
  \end{scope}
\end{tikzpicture}
\end{center}

Comparación con los autovectores de \(A\): \(\lambda=5\) con \((0,1)\) y \(\lambda=3\)
con \((1,-2)\) (verificar: \(A(1,-2)^T=(3,-6)^T\)). No son perpendiculares entre sí y no
coinciden con \(v_1,v_2\). Mismo producto \(3\cdot5=15\), pero ejes distintos: la
diagonalización mira las direcciones que caen sobre sí mismas; la SVD mira las
direcciones perpendiculares que siguen perpendiculares. Dos casos límite: la rotación de
\(90^\circ\) tiene \(R^TR=I\), valores singulares \(1,1\), la circunferencia sigue siendo
circunferencia (``no estira nada, sólo gira''); la simétrica
\(\begin{pmatrix}2&1\\1&2\end{pmatrix}\) del capítulo~\ref{cap:cambio-base} (sección~\ref{sec:cambio-base}) tiene valores singulares
\(3,1\) en las direcciones \((1,1),(1,-1)\): acá autovectores y vectores singulares
coinciden, \(U=V\) y \(\sigma_i=|\lambda_i|\). El resumen 2025 de la clase 10 dice que en
\(\R^2\) las simétricas ``no son un caso de coincidencia sino El caso''.

\paragraph{El cuaderno del profesor: una \texorpdfstring{$2\times3$}{2x3} analítica.}
\file{LaSVD.wxmx} (apéndice~\ref{ap:notebooks}) hace en wxMaxima la SVD de
\(A=\begin{pmatrix}1&2&4\\-3&-2&3\end{pmatrix}\): calcula \(AA^T\) y \(A^TA\), el
polinomio característico, los autovectores, verifica que son ortogonales de a pares,
normaliza \(U\) y \(V^T\), arma \(\Sigma\) con las raíces de los autovalores y comprueba
\(U\Sigma V^T=A\). Rehecho acá a mano y con NumPy:
\[
  AA^T=\begin{pmatrix}21&5\\5&22\end{pmatrix},\qquad
  \mu^2-43\mu+437=0\ \Rightarrow\ \mu=\frac{43\pm\sqrt{101}}2
  \ \Rightarrow\ \sigma_1=\sqrt{\mu_1}\approx5{,}150,\ \ \sigma_2=\sqrt{\mu_2}\approx4{,}059 .
\]
Autovectores de \(AA^T\) (fila \((21-\mu,\ 5)\) de \(AA^T-\mu I\), ortogonal al autovector,
así que \(u\propto(5,\ \mu-21)\)): \(u_1\approx(0{,}671,\ 0{,}741)\),
\(u_2\approx(0{,}741,\ -0{,}671)\). La base de entrada se obtiene al revés que en el
ejemplo anterior, \(v_i=A^Tu_i/\sigma_i\):
\(v_1\approx(-0{,}302,\ -0{,}027,\ 0{,}953)\), \(v_2\approx(0{,}679,\ 0{,}696,\ 0{,}235)\),
y el tercer vector de \(V\) es el núcleo de \(A\), el producto cruz de sus dos filas
(sección~\ref{sec:nucleo}): \((14,-15,4)/\sqrt{437}\approx(0{,}670,\ -0{,}718,\ 0{,}191)\).
Chequeo: \(Av_1=\sigma_1u_1\) da \((3{,}456,\ 3{,}819)\) por los dos caminos, y
\(A^TA=\begin{pmatrix}10&8&-5\\8&8&2\\-5&2&25\end{pmatrix}\) tiene autovalores
\(26{,}525\), \(16{,}475\) y \(0\). Es el caso rectangular: \(U\) de \(2\times2\),
\(\Sigma=\begin{pmatrix}\sigma_1&0&0\\0&\sigma_2&0\end{pmatrix}\) de \(2\times3\), \(V\) de
\(3\times3\), ``una base a la
entrada que no sale'': la esfera unidad de \(\R^3\) se aplasta a una elipse de
\(\R^2\) y la dirección \(v_3\) cae al origen. El cuaderno del profesor hace exactamente esto: \texttt{solve} del polinomio de
\(A^TA\), \texttt{eigenvectors} de las dos matrices, \texttt{ratsimp} de los productos
internos de a pares para ver los ceros, normalización, y \texttt{float(U.\(\Sigma\).Vt)}
que devuelve \(A\). Las raíces son irracionales: es lo normal, no un error
(caja ``Sobre los números lindos'' de la sección~\ref{sec:autovalores}). En NumPy,
\texttt{np.linalg.svd(A)} devuelve \texttt{U, s, Vt}; los signos de las columnas pueden
venir cambiados, sólo importa la dirección. La tarea que dejó el profesor el 14/10 es
reproducir esto en GeoGebra con cualquier \(2\times3\) o \(3\times2\).

\section{En profundidad}


Aplicación al terreno: véase el complemento \ref{cap:dem}.
